% dimensional_selection.tex
% Section 9 (NEW in v0.2) -- the headline corollary: a minimum
% probability quantum δp_min forces d_max = 1/δp_min - 1.  At δp_min
% = 1/4, d_max = 3.  This is the candidate mechanism that selects
% the universe's dimension rather than permitting it.

\label{sec:dimensional_selection}

\placeholder{Section body -- write during Phase~2.  State and prove
the dimensional-selection corollary; spell out the
$\delta p_\min = 1/4 \Rightarrow d_\max = 3$ specialisation; document
the cross-reference to \texttt{dcl-delta-p-min} for the empirical
investigation of $\delta p_\min$'s value.}

\subsection{Setup: A Minimum Probability Quantum}
\label{subsec:dp_min_setup}

\placeholder{Take $\delta p_\min > 0$ as a hypothesis: the A=1
framework admits no probability per site smaller than $\delta p_\min$.
The hypothesis can be realised in (at least) two operationally
distinct ways:

\begin{itemize}
\item A \emph{clamp} on continuous amplitudes: $|psi(x)|^2 \geq
   \delta p_\min$ when non-zero.  Implementation: the
   \texttt{prob\_floor} parameter on \texttt{dcl\_core.core}
   (planned \texttt{dcl-core v0.2.0}).
\item An \emph{intrinsic granularity}: probability is represented as
   a multiple of $\delta p_\min$.  Implementation: the integer-token
   engine \texttt{dcl\_core.core3d} with token budget
   $N = 1/\delta p_\min$.
\end{itemize}

The cross-engine equivalence question -- whether these two
realisations produce the same physics -- is the open question of
Phase~1 in \texttt{dcl-delta-p-min}.  This section's result is
operationally engine-agnostic: it depends only on the *existence* of
a $\delta p_\min > 0$, not on its realisation.}

\subsection{The Dimensional-Selection Corollary}
\label{subsec:dp_min_dmax_corollary}

\begin{theorem}[Dimensional selection from $\delta p_\min$]
\label{thm:dp_min_dmax}
If a minimum probability quantum $\delta p_\min > 0$ exists in the
A=1 framework, then the diamond progression $\Lambda_d$ can sustain
the $\mathcal{A} = 1$ unity axiom only for
\begin{equation*}
d \leq d_{\max} = \frac{1}{\delta p_\min} - 1.
\end{equation*}
\end{theorem}

\begin{proof}
\placeholder{Proof: at any site $x$ of $\Lambda_d$, the unity
constraint $\mathcal{A} = 1$ requires the per-site probability mass
to distribute over $x$ together with all the configurations
participating in the joint-tick rule.  By
Theorem~\ref{thm:coord_2d}, the site has $2d$ neighbours; combined
with the joint-tick-rule's structure, $d + 1$ distinct probability
slots receive non-zero mass.  Each slot must carry probability at
least $\delta p_\min$, so
$(d+1) \cdot \delta p_\min \leq 1$, giving
$d \leq 1/\delta p_\min - 1$.  Spell out the joint-tick-rule
participation count once the construction is fixed in
Section~\ref{sec:lattice_n_dim}.}
\end{proof}

\paragraph{Note on the proof structure.}
\todomark{motivating-example: the exact slot-count ``$d+1$'' depends
on a specific reading of the joint-tick rule that pulls from Paper~I
plus the diamond-progression structure of
Section~\ref{sec:lattice_n_dim}.  Phase~2 work must verify that
this count is precise rather than a sketch; alternative readings
(e.g.\ $2d+1$ slots) give different $d_\max$ formulae.  The
formula $d_\max = 1/\delta p_\min - 1$ is the user's stated result
from `raw\_thoughts.md` 2026-05-21 \texttt{[result-2]}; this paper
re-derives and confirms it.}

\subsection{Specialisation: $\delta p_\min = 1/4 \Rightarrow d_{\max} = 3$}
\label{subsec:dp_min_dmax_3}

\begin{corollary}
\label{cor:dmax_3}
If $\delta p_\min = 1/4$, then $d_{\max} = 3$.
\end{corollary}

\begin{proof}
Direct substitution into Theorem~\ref{thm:dp_min_dmax}:
$d_{\max} = 1/(1/4) - 1 = 4 - 1 = 3$.
\end{proof}

\paragraph{Significance.}
This corollary is the candidate mechanism that \emph{forces} the
universe to be three-dimensional rather than permitting it.  If the
A=1 framework forces $\delta p_\min = 1/4$ (via a token-conservation
constraint or analogous structural argument), then it forces $d =
3$.  The framework is no longer compatible with our observed
universe by accident -- it is compatible because the unique
admissible $d$ is $3$.

\paragraph{Connection to Paper~I.}
Paper~I established the bipartite octahedral lattice as the
substrate of physics in $(3+1)$ dimensions.  At the time of
Paper~I, the specific choice of $d = 3$ was \emph{motivated} by
observation -- physics is 3D, so the lattice is 3D.  This
corollary inverts the relationship: given $\delta p_\min$, the
framework forces $d = 3$.  Observational compatibility becomes
a structural consequence rather than a motivating constraint.

\subsection{The Empirical Question: Is $\delta p_\min = 1/4$?}
\label{subsec:dp_min_value_question}

\placeholder{The corollary above is mathematical: it derives
$d_\max$ from $\delta p_\min$ given the diamond progression's
structure.  Whether the framework forces $\delta p_\min$ to take a
specific value (and which value) is a separate question -- empirical
and / or structural-from-token-conservation.

The companion repository \texttt{dcl-delta-p-min} runs the
investigation.  Briefly: a 4-cell numerical grid (two engines
crossed with two $\delta p_\min$ regimes) measures how the
Arnold-tongue width and hydrogen Bohr-radius recovery depend on
$\delta p_\min$; the largest $\delta p_\min$ at which Paper~I's
PASS audit rows remain recoverable is a framework-internal upper
bound.  The planned discrete-probability paper plausibly pins
$\delta p_\min$'s natural value via a token-conservation argument
(Outcome~D of \texttt{dcl-delta-p-min}'s Phase~4 stance decision).
At $\delta p_\min = 1/4$, Corollary~\ref{cor:dmax_3} forces $d_\max
= 3$ exactly; at other values, the dimensional ceiling shifts.

For this paper, the result stands as a \emph{conditional theorem}:
if $\delta p_\min$ exists and equals $1/4$, then $d_\max = 3$.
The conditional is non-empty -- the diamond progression admits
$\delta p_\min$ structurally, and the value $1/4$ is the empirical
target of \texttt{dcl-delta-p-min}.}

\subsection{Implications}
\label{subsec:dimensional_selection_implications}

\begin{itemize}
\item \emph{The Hilbert-Sixth capstone}
      (\texttt{dcl-paper-XX-hilbert-sixth}, planned) consumes this
      corollary as part of its universality argument: the
      axiomatisation forces $d = 3$.
\item \emph{The planned discrete-probability paper} fixes
      $\delta p_\min$'s value (Outcome~D of the
      \texttt{dcl-delta-p-min} investigation); this paper's
      corollary then becomes a parameter-free statement.
\item \emph{The aspirational F4 sequel} (`A=1 forces linear
      coordination', tracked in
      \texttt{physics-research/notes/raw\_thoughts.md}) asks
      whether the diamond progression is the \emph{only} bipartite
      lattice family admitting A=1 non-degenerately.  If F4 lands,
      Corollary~\ref{cor:dmax_3} sharpens further: not just
      ``$d_\max = 3$ for the diamond progression'' but ``$d_\max
      = 3$ for any lattice family admitting A=1''.
\end{itemize}
